\section{A cost-dependent contraction window}
The critical-ratio theorem implies that ``when should contraction start?'' has no architecture-free answer even in the exact rank-one phase model. The statistical geometry determines when useful target structure becomes observable, while the realized wide-to-compact cost ratio determines whether waiting for that structure is economically worthwhile.

\begin{definition}[Compute-feasible contraction set]\label{def:compute-window}
Fix a compact continuation horizon $H$ and a realized wide-to-compact cost ratio $\rho=C_W/C_1\ge1$. In the rank-one phase model define
\[
\mathcal T_\rho(H)
:=
\{t>0:\rho<\rho_\star(t,H)\}.
\]
When the set is nonempty, define its earliest and latest feasible times
\[
t_{\rm on}(\rho,H):=\inf\mathcal T_\rho(H),
\qquad
 t_{\rm off}(\rho,H):=\sup\mathcal T_\rho(H).
\]
No connectedness or monotonicity of $\mathcal T_\rho(H)$ is assumed.
\end{definition}

\begin{corollary}[Exact cost-dependent onset]\label{cor:cost-dependent-onset}
Under the assumptions of Theorem~\ref{thm:rankone-critical-compute}, a wide-then-rank-one contraction performed at time $t$ beats the matched rank-one-from-start path at equal compute if and only if
\[
t\in\mathcal T_\rho(H).
\]
If a separate statistical safety condition requires $t\ge\tau_{\rm safe}$, then the earliest oracle contraction time that is both statistically safe and compute-beneficial is
\[
\boxed{
 t_{\rm contract}(\rho,H)
=
\inf\bigl(\mathcal T_\rho(H)\cap[\tau_{\rm safe},\infty)\bigr),}
\]
with value $+\infty$ if the intersection is empty.
\end{corollary}

\begin{proof}
The first statement is the definition of $\mathcal T_\rho(H)$ combined with Theorem~\ref{thm:rankone-critical-compute}. Adding the independent safety requirement restricts the admissible set to the displayed intersection; its infimum is therefore the earliest feasible contraction time.
\end{proof}

\begin{remark}[There can be a latest useful time]
Statistical safety generally improves as target overlap increases and the nuisance bulk contracts. Compute value need not improve forever. Once the target-bearing direction is already resolved, additional wide training can add little discovery value while continuing to pay the wide cost. Hence $\rho_\star(t,H)$ may eventually fall, and $\mathcal T_\rho(H)$ may have a finite upper edge $t_{\rm off}$. This possibility is important algorithmically: ``wait longer to be safer'' can become economically wrong after the discovery phase has ended.
\end{remark}

\begin{remark}[Three distinct clocks]
The theory therefore contains three clocks that should not be collapsed into one heuristic threshold:
\[
\boxed{
\text{spectral resolution time}
\quad\neq\quad
\text{safe contraction time}
\quad\neq\quad
\text{compute-optimal contraction time}.}
\]
The first concerns identification of target-bearing geometry, the second bounds statistical damage from deletion, and the third compares the discovery value of remaining wide with the compute opportunity cost of doing so. A practical controller must satisfy all three constraints simultaneously.
\end{remark}
